\begin{itemize}
\item \textbf{NORMAL.} \emph{Numérica:} Sin ruido y con pocos nodos, el polinomio de grado 12 es muy preciso (error 3.323$\times10^{-7}$ en los ocultos) y supera al spline (0.00117), porque sin t es analítica y sus derivadas están acotadas. Lagrange y Newton difieren sólo en 1.549$\times10^{-14}$; Newton permite agregar nodos sin recalcular todo y evaluar con Horner. \emph{Problema:} Para el sensor, la reconstrucción de los 4 registros perdidos es confiable con Newton (error máximo 3.323$\times10^{-7}$ u). El cruce del umbral se estima en t = 0.524165 s frente a 0.523599 s reales.
\item \textbf{ESTRÉS.} \emph{Numérica:} Con ruido $\sigma$ = 0.002 y grado 24, el polinomio amplifica el ruido hasta 39.5 ($\Lambda$ = 1.86$\times10^{5}$), mientras el spline lo deja en 0.00543. En los registros ocultos el spline logra 0.00533 frente a 4.31 del polinomio. El sistema del spline converge con CG, Gauss-Seidel y Jacobi por ser diagonal dominante. \emph{Problema:} Para el sensor, la reconstrucción de los 6 registros perdidos es confiable con Spline cúbico natural (error máximo 0.00533 u). Con ruido, un polinomio de grado alto no debe usarse para rellenar huecos: convierte pequeñas perturbaciones en grandes errores; el spline es la opción robusta. El cruce del umbral se estima en t = 1.57021 s frente a 1.5708 s reales.
\item \textbf{MAL COND..} \emph{Numérica:} Con 21 nodos uniformes, $\Lambda$ = 1.099$\times10^{4}$ y $\kappa_{2}$(Vandermonde) = 8.314$\times10^{8}$: el polinomio interpolante oscila cerca de los extremos (error 59.8) aunque Lagrange y Newton coinciden hasta 8.095$\times10^{-11}$. No es un fallo de los algoritmos: con nodos de Chebyshev ($\Lambda$ = 2.9) el mismo grado da 0.0153, y el spline, de grado bajo por tramos, da 0.00318. \emph{Problema:} Para el sensor, la reconstrucción de los 20 registros perdidos es confiable con Spline cúbico natural (error máximo 0.00317 u). Si el muestreo es uniforme y la señal tiene un pulso, sólo el spline (o remuestrear en nodos de Chebyshev) da una reconstrucción útil cerca de los bordes. El cruce del umbral se estima en t = -0.2 s frente a -0.2 s reales.
\end{itemize}
